\documentclass{ximera}

\begin{document}
	\author{Stitz-Zeager}
	\xmtitle{Exercises for Polar Complex}{}

\mfpicnumber{1} \opengraphsfile{ExercisesforPolarComplex} % mfpic settings added 

\begin{question}
In Exercises \ref{polarcompbasicfirst} - \ref{polarcompbasiclast}, find a polar representation for the complex number $z$.  Identify $\text{Re}(z)$, $\text{Im}(z)$, $|z|$, $\text{arg}(z)$ and $\text{Arg}(z)$.

\begin{problem} \label{polarcompbasicfirst}
$z = 9 + 9i$ 

\begin{solution}
$z =  9\sqrt{2}\text{cis}\left(\frac{\pi}{4}\right)$

$\text{Re}(z) = 9$

$\text{Im}(z) = 9$

$|z| = 9\sqrt{2}$

$\text{arg}(z) = \left\{\frac{\pi}{4} + 2\pi k \, | \, \text{$k$ is an integer} \right\}$

$\text{Arg}(z) = \frac{\pi}{4}$

\end{solution}
\end{problem}

\begin{problem}
$z = 5 + 5i\sqrt{3}$

$z =\answer{10}\text{cis}\left(\answer{\frac{\pi}{3}}\right)$  (Choose $\theta$ so $0 \leq \theta \leq 2\pi$.)

$\text{Re}(z) = \answer{5}$

$\text{Im}(z) = \answer{5\sqrt{3}}$

$|z| = \answer{10}$

$\text{arg}(z) = \left\{\answer{\frac{\pi}{3}} + 2\pi k \, | \, \text{$k$ is an integer} \right\}$  (Choose $\theta$ so $0 \leq \theta \leq 2\pi$.)

$\text{Arg}(z) = \answer{\frac{\pi}{3}}$


\end{problem}

\begin{problem}
$z = 6i$

\begin{solution}
$z = 6\text{cis}\left(\frac{\pi}{2}\right)$

$\text{Re}(z) = 0$

$\text{Im}(z) = 6$

$|z| = 6$

$\text{arg}(z) = \left\{\frac{\pi}{2} + 2\pi k \, | \, \text{$k$ is an integer} \right\}$

$\text{Arg}(z) = \frac{\pi}{2}$

\end{solution}
\end{problem}

\begin{problem}
$z = -3\sqrt{2} + 3i\sqrt{2}$

$z = \answer{6}\text{cis}\left(\answer{\frac{3\pi}{4}}\right)$ (Choose $\theta$ so $0 \leq \theta \leq 2\pi$.)

$\text{Re}(z) = \answer{-3\sqrt{2}}$

$\text{Im}(z) = \answer{3\sqrt{2}}$

$|z| = \answer{6}$

$\text{arg}(z) = \left\{\answer{\frac{3\pi}{4}} + 2\pi k \, | \, \text{$k$ is an integer}  \right\}$ (Choose $\theta$ so $0 \leq \theta \leq 2\pi$.)

$\text{Arg}(z) = \answer{\frac{3\pi}{4}}$



\end{problem}

\begin{problem}
$z = -6\sqrt{3} + 6i$

\begin{solution}

$z =  12\text{cis}\left(\frac{5\pi}{6}\right)$

$\text{Re}(z) = -6\sqrt{3}$

$\text{Im}(z) =6$

$|z| = 12$

$\text{arg}(z) = \left\{\frac{5\pi}{6} + 2\pi k \, | \, \text{$k$ is an integer} \right\}$

$\text{Arg}(z) = \frac{5\pi}{6}$

\end{solution}
\end{problem}

\begin{problem}
$z = -2$ 

$z = \answer{2}\text{cis}\left(\answer{\pi}\right)$ (Choose $\theta$ so $0 \leq \theta \leq 2\pi$.)

$\text{Re}(z) = \answer{-2}$

$\text{Im}(z) =\answer{0}$

$|z| = \answer{2}$

$\text{arg}(z) = \left\{ \answer{\pi} + 2\pi k = (2k+1)\pi \, | \, \text{$k$ is an integer} \right\}$ (Choose $\theta$ so $0 \leq \theta \leq 2\pi$.)

$\text{Arg}(z) = \answer{\pi}$


\end{problem}

\begin{problem}
$z = -\frac{\sqrt{3}}{2} - \frac{1}{2}i$

\begin{solution}

$z = \text{cis}\left(\frac{7\pi}{6}\right)$

$\text{Re}(z) = -\frac{\sqrt{3}}{2}$

$\text{Im}(z) = -\frac{1}{2}$

$|z| = 1$

$\text{arg}(z) = \left\{\frac{7\pi}{6} + 2\pi k \, | \, \text{$k$ is an integer} \right\}$

$\text{Arg}(z) = -\frac{5\pi}{6}$

\end{solution}

\end{problem}

\begin{problem}
$z = -3-3i$ 

$z = \answer{3\sqrt{2}}\text{cis}\left(\answer{\frac{5\pi}{4}}\right)$ (Choose $\theta$ so $0 \leq \theta \leq 2\pi$.)

$\text{Re}(z) =\answer{ -3}$

$\text{Im}(z) =\answer{-3}$

$|z| = \answer{3\sqrt{2}}$

$\text{arg}(z) = \left\{\answer{\frac{5\pi}{4}} + 2\pi k \, | \, \text{$k$ is an integer} \right\}$ (Choose $\theta$ so $0 \leq \theta \leq 2\pi$.)

$\text{Arg}(z) = \answer{-\frac{3\pi}{4}}$

\end{problem}

\begin{problem}
 $z = -5i$

 \begin{solution}
$z = 5\text{cis}\left(\frac{3\pi}{2}\right)$

$\text{Re}(z) = 0$

$\text{Im}(z) = -5$

$|z| = 5$

$\text{arg}(z) = \left\{\frac{3\pi}{2} + 2\pi k \, | \, \text{$k$ is an integer} \right\}$

$\text{Arg}(z) = -\frac{\pi}{2}$

\end{solution}
\end{problem}

\begin{problem}
 $z = 2\sqrt{2} - 2i\sqrt{2}$

 $z = \answer{4}\text{cis}\left(\answer{\frac{7\pi}{4}}\right)$  (Choose $\theta$ so $0 \leq \theta \leq 2\pi$.)
 
 $\text{Re}(z) = \answer{2\sqrt{2}}$
 
 $\text{Im}(z) = \answer{-2\sqrt{2}}$
 
 $|z| = \answer{4}$

$\text{arg}(z) = \left\{\answer{\frac{7\pi}{4}} + 2\pi k \, | \, \text{$k$ is an integer} \right\}$ (Choose $\theta$ so $0 \leq \theta \leq 2\pi$.)

$\text{Arg}(z) = \answer{-\frac{\pi}{4}}$

\end{problem}

\begin{problem}
$z = 6$

\begin{solution}
$z = 6\text{cis}\left(0\right)$

$\text{Re}(z) = 6$

$\text{Im}(z) = 0$

$|z| = 6$

$\text{arg}(z) = \left\{0+2\pi k  = 2 \pi k\, | \, \text{$k$ is an integer} \right\}$ 

$\text{Arg}(z) =0$

\end{solution}
\end{problem}

\begin{problem}
$z = i\sqrt[3]{7}$

 $z = \answer{\sqrt[3]{7}}\text{cis}\left(\answer{\frac{\pi}{2}}\right)$ (Choose $\theta$ so $0 \leq \theta \leq 2\pi$.)
 
 $\text{Re}(z) =\answer{0}$
 
 $\text{Im}(z) = \answer{\sqrt[3]{7}}$
 
 $|z| = \answer{\sqrt[3]{7}}$

$\text{arg}(z) = \left\{\answer{\frac{\pi}{2}} + 2\pi k \, | \, \text{$k$ is an integer} \right\}$  (Choose $\theta$ so $0 \leq \theta \leq 2\pi$.)

$\text{Arg}(z) = \answer{\frac{\pi}{2}}$
\end{problem}

\begin{problem}
$z = 3 + 4i$

\begin{solution}

$z = 5\text{cis}\left(\arctan\left(\frac{4}{3}\right)\right)$

$\text{Re}(z) = 3$

$\text{Im}(z) = 4$

$|z| = 5$

$\text{arg}(z) = \left\{\arctan\left(\frac{4}{3}\right) + 2\pi k \, | \, \text{$k$ is an integer} \right\}$

$\text{Arg}(z) =\arctan\left(\frac{4}{3}\right) $

\end{solution}
\end{problem}

\begin{problem}
$z = \sqrt{2} + i$

$z = \answer{\sqrt{3}}\text{cis}\left(\answer{\arctan\left(\frac{\sqrt{2}}{2}\right)}\right)$ (Choose $\theta$ so $0 \leq \theta \leq 2\pi$.)

$\text{Re}(z) = \answer{\sqrt{2}}$

$\text{Im}(z) = \answer{1}$

$|z| = \answer{\sqrt{3}}$

$\text{arg}(z) = \left\{\answer{\arctan\left(\frac{\sqrt{2}}{2}\right)} + 2\pi k \, | \, \text{$k$ is an integer} \right\}$  (Choose $\theta$ so $0 \leq \theta \leq 2\pi$.)

$\text{Arg}(z) =\answer{\arctan\left(\frac{\sqrt{2}}{2}\right)} $

\end{problem}

\begin{problem}
$z = -7 + 24i$

\begin{solution}
$z = 25\text{cis}\left(\pi - \arctan\left(\frac{24}{7}\right)\right)$

$\text{Re}(z) = -7$

$\text{Im}(z) = 24$

$|z| = 25$

$\text{arg}(z) = \left\{\pi - \arctan\left(\frac{24}{7}\right) + 2\pi k \, | \, \text{$k$ is an integer} \right\}$

$\text{Arg}(z) =\pi - \arctan\left(\frac{24}{7}\right) $

\end{solution}
\end{problem}

\begin{problem}
$z = -2+6i$

$z = \answer{2\sqrt{10}}\text{cis}\left(\answer{\pi - \arctan\left(3\right)}\right)$ (Choose $\theta$ so $0 \leq \theta \leq 2\pi$.)

$\text{Re}(z) = \answer{-2}$

$\text{Im}(z) = \answer{6}$

$|z| =\answer{2\sqrt{10}}$

$\text{arg}(z) = \left\{ \answer{\pi - \arctan\left(3\right)} + 2\pi k \, | \, \text{$k$ is an integer} \right\}$ (Choose $\theta$ so $0 \leq \theta \leq 2\pi$.)

$\text{Arg}(z) = \answer{\pi - \arctan\left(3\right)} $.

\end{problem}

\begin{problem}
$z = -12-5i$

\begin{solution}

$z =13\text{cis}\left(\pi + \arctan\left(\frac{5}{12}\right)\right)$

$\text{Re}(z) = -12$

$\text{Im}(z) = -5$

$|z| = 13$

$\text{arg}(z) = \left\{\pi +\arctan\left(\frac{5}{12}\right) + 2\pi k \, | \, \text{$k$ is an integer} \right\}$

$\text{Arg}(z) =  \arctan\left(\frac{5}{12}\right) -\pi $.

\end{solution}

\end{problem}

\begin{problem}
$z = -5-2i$

$z = \answer{\sqrt{29}}\text{cis}\left(\answer{\pi + \arctan\left(\frac{2}{5}\right)}\right)$ (Choose $\theta$ so $0 \leq \theta \leq 2\pi$.)

$\text{Re}(z) = \answer{-5}$

$\text{Im}(z) = \answer{-2}$

$|z| = \answer{\sqrt{29}}$

$\text{arg}(z) = \left\{\answer{\pi +\arctan\left(\frac{2}{5}\right)} + 2\pi k \, | \, \text{$k$ is an integer} \right\}$

$\text{Arg}(z) =  \answer{\arctan\left(\frac{2}{5}\right) -\pi} $.

\end{problem}

\begin{problem}
$z = 4-2i$

\begin{solution}
$z = 2\sqrt{5}\text{cis}\left(\arctan\left(-\frac{1}{2}\right)\right)$

$\text{Re}(z) =4$

$\text{Im}(z) = -2$

$|z| = 2\sqrt{5}$

$\text{arg}(z) = \left\{\arctan\left(-\frac{1}{2}\right) + 2\pi k \, | \, \text{$k$ is an integer} \right\}$

$\text{Arg}(z) = \arctan\left(-\frac{1}{2}\right) = -\arctan\left(\frac{1}{2}\right) $.
\end{solution}
\end{problem}

\begin{problem}\label{polarcompbasiclast}
$z = 1-3i$ 

$z = \answer{\sqrt{10}}\text{cis}\left(\answer{\arctan\left(-3\right)}\right)$ (Choose $\theta$ so $-\pi < \theta \leq \pi$.)

$\text{Re}(z) =\answer{1}$

$\text{Im}(z) = \answer{-3}$

$|z| = \answer{\sqrt{10}}$

$\text{arg}(z) = \left\{ \answer{\arctan\left(-3\right)} + 2\pi k \, | \, \text{$k$ is an integer} \right\}$ (Choose $\theta$ so $-\pi < \theta \leq \pi$.)

$\text{Arg}(z) =  \answer{\arctan\left(-3\right)}$.

\end{problem}

\end{question}


\begin{question}
In Exercises \ref{rectcompfirst} - \ref{rectcomplast}, find the rectangular form of the given complex number.  Use whatever identities are necessary to find the exact values.
\begin{problem} \label{rectcompfirst}
 $z = 6\text{cis}(0)$  

 \begin{solution}
$z  = 6$
 \end{solution}
\end{problem}

\begin{problem}
$z = 2\text{cis}\left(\frac{\pi}{6}\right)$

$z  = \answer{\sqrt{3}} + \answer{1}i$

\end{problem}

\begin{problem}
$z = 7\sqrt{2}\text{cis}\left(\frac{\pi}{4}\right)$

\begin{solution}
$z = 7+7i$
\end{solution}

\end{problem}

\begin{problem}
$z = 3\text{cis}\left(\frac{\pi}{2}\right)$

 $z = \answer{0} + \answer{3}i$ 
\end{problem}

\begin{problem}
$z = 4\text{cis}\left(\frac{2\pi}{3}\right)$

\begin{solution}

$z = -2+2i\sqrt{3}$
\end{solution}
\end{problem}

\begin{problem}
$z = \sqrt{6}\text{cis}\left(\frac{3\pi}{4}\right)$ 

$z = \answer{-\sqrt{3}} + \answer{\sqrt{3}} \, i$ 

\end{problem}

\begin{problem}
$z = 9\text{cis}\left(\pi\right)$

\begin{solution}
$z = -9$
\end{solution}

\end{problem}

\begin{problem}
$z = 3\text{cis}\left(\frac{4\pi}{3}\right)$

$z = \answer{-\frac{3}{2}} + \answer{-\frac{3\sqrt{3}}{2}} i$

\end{problem}

\begin{problem}
$z = 7\text{cis}\left(-\frac{3\pi}{4}\right)$

\begin{solution}
$z = -\frac{7\sqrt{2}}{2} - \frac{7\sqrt{2}}{2}i$ 
\end{solution}
\end{problem}

\begin{problem}
$z = \sqrt{13}\text{cis}\left(\frac{3\pi}{2}\right)$ 

$z = \answer{-\sqrt{13}} \, i$ 

\end{problem}

\begin{problem}
$z = \frac{1}{2}\text{cis}\left(\frac{7\pi}{4}\right)$ 

\begin{solution}
 $z = \frac{\sqrt{2}}{4} - i\frac{\sqrt{2}}{4}$ 
\end{solution}
\end{problem}

\begin{problem}
$z = 12\text{cis}\left(-\frac{\pi}{3}\right)$ 

 $z = \answer{6} + \answer{-6\sqrt{3}} \, i$ 
\end{problem}

\begin{problem}
$z = 8\text{cis}\left(\frac{\pi}{12}\right)$

\begin{solution}
$z = 4\sqrt{2+\sqrt{3}}+4i\sqrt{2-\sqrt{3}}$ 
\end{solution}
\end{problem}

\begin{problem}
$z = 2\text{cis}\left(\frac{7\pi}{8}\right)$

$z = \answer{-\sqrt{2 + \sqrt{2}}} + \answer{\sqrt{2 - \sqrt{2}}} \, i $ 
\end{problem}

\begin{problem}
$z = 5\text{cis}\left(\arctan\left(\frac{4}{3}\right)\right)$

\begin{solution}
$z  = 3 + 4i$ 
\end{solution}
\end{problem}

\begin{problem}
$z = \sqrt{10}\text{cis}\left(\arctan\left(\frac{1}{3}\right)\right)$ 

$z = \answer{3}+\answer{1} \, i$ 

\end{problem}

\begin{problem}
$z = 15\text{cis}\left(\arctan\left(-2\right)\right)$

\begin{solution}
$z = 3\sqrt{5} -6i\sqrt{5}$ 
\end{solution}
\end{problem}

\begin{problem}
$z=  \sqrt{3}\text{cis}\left(\arctan\left(-\sqrt{2}\right)\right)$

 $z = \answer{1}+ \answer{-\sqrt{2} \, i}$

\end{problem}

\begin{problem}
$z = 50\text{cis}\left(\pi-\arctan\left(\frac{7}{24}\right)\right)$

\begin{solution}
 $z = -48 + 14i$ 
\end{solution}
\end{problem}

\begin{problem} \label{rectcomplast}
$z = \frac{1}{2}\text{cis}\left(\pi+\arctan\left(\frac{5}{12}\right)\right)$ 

$z = \answer{-\frac{6}{13}} +  \answer{-\frac{5}{26}} i$
\end{problem}

\end{question}

\begin{question}
For Exercises \ref{polarcomparithfirst} - \ref{polarcomparithlast}, use $z = -\frac{3\sqrt{3}}{2} + \frac{3}{2}i$ and $w = 3\sqrt{2} - 3i\sqrt{2}$ to compute the quantity.  Express your answers in polar form using the principal argument.

\begin{problem} \label{polarcomparithfirst}
 $zw$

 \begin{solution}
$zw = 18\text{cis}\left(\frac{7\pi}{12}\right)$
 \end{solution}
\end{problem}

\begin{problem}
$\frac{z}{w}$

\begin{solution}
$\frac{z}{w} = \frac{1}{2}\text{cis}\left(-\frac{11\pi}{12}\right)$
\end{solution}
\end{problem}

\begin{problem}
$\frac{w}{z}$

$\frac{w}{z} = \answer{2}\text{cis}\left(\answer{\frac{11\pi}{12}}\right)$

\end{problem}

\begin{problem}
$z^{4}$

\begin{solution}
$z^{4} = 81\text{cis}\left(-\frac{2\pi}{3}\right)$
\end{solution}
\end{problem}

\begin{problem}
$w^{3}$

$w^{3} = \answer{216}\text{cis}\left(\answer{-\frac{3\pi}{4}}\right)$

\end{problem}

\begin{problem}
$z^{5}w^{2}$

\begin{solution}
$z^{5}w^{2} = 8748\text{cis}\left(-\frac{\pi}{3}\right)$
\end{solution}
\end{problem}

\begin{problem}
$z^{3}w^{2}$

$z^3w^2 = \answer{972} \text{cis}(\answer{0})$
\end{problem}

\begin{problem}
$\frac{z^{2}}{w}$

\begin{solution}
 $\frac{z^2}{w} =\frac{3}{2}\text{cis}\left(-\frac{\pi}{12}\right)$
\end{solution}
\end{problem}

\begin{problem}
$\frac{w}{z^2}$

 $\frac{w}{z^2} =\answer{\frac{2}{3}}\text{cis}\left(\answer{\frac{\pi}{12}}\right)$
\end{problem}

\begin{problem}
$\frac{z^3}{w^2}$

\begin{solution}
 $\frac{z^3}{w^2} =\frac{3}{4}\text{cis}(\pi)$
\end{solution}
\end{problem}

\begin{problem}
$\frac{w^2}{z^3}$

$\frac{w^2}{z^3} =\answer{\frac{4}{3}}\text{cis}(\answer{\pi})$
\end{problem}

\begin{problem}  \label{polarcomparithlast}
$\left(\frac{w}{z}\right)^6$ 

\begin{solution}
 $\left(\frac{w}{z}\right)^6 =64\text{cis}\left(-\frac{\pi}{2} \right)$
\end{solution}
 
\end{problem}

\end{question}

\begin{question}
In Exercises \ref{demoivrefirst} - \ref{demoivrelast}, use DeMoivre's Theorem to find the indicated power of the given complex number.  Express your final answers in rectangular form.

\begin{problem}\label{demoivrefirst}
$\left(-2 + 2i\sqrt{3}\right)^3$ 

\begin{solution}
$\left(-2 + 2i\sqrt{3}\right)^3 = 64$
\end{solution}
\end{problem}

\begin{problem}
$(-\sqrt{3} - i)^3$ 

 $(-\sqrt{3} - i)^3 = \answer{0} + \answer{-8} \, i$
\end{problem}

\begin{problem}
$(-3+3i)^{4}$

\begin{solution}
$(-3+3i)^{4}=-324$
\end{solution}
\end{problem}

\begin{problem}
$(\sqrt{3} + i)^4$

$(\sqrt{3} + i)^4 =\answer{-8} + \answer{8\sqrt{3}} \, i$
\end{problem}

\begin{problem}
$\left(\frac{5}{2} + \frac{5}{2} i\right)^3$ 

\begin{solution}
$\left(\frac{5}{2} + \frac{5}{2} i\right)^3=-\frac{125}{4}+\frac{125}{4} i$ 
\end{solution}
\end{problem}

\begin{problem}
$\left(-\frac{1}{2} - \frac{\sqrt{3}}{2} i\right)^{6}$

$\left(-\frac{1}{2} - \frac{i \sqrt{3}}{2}\right)^{6}=\answer{1} + \answer{0} i$
\end{problem}

\begin{problem}
$\left(\frac{3}{2} - \frac{3}{2} i\right)^3$ 

\begin{solution}
$\left(\frac{3}{2} - \frac{3}{2} i\right)^3=-\frac{27}{4}-\frac{27}{4} i$ 
\end{solution}
\end{problem}

\begin{problem}
$\left(\frac{\sqrt{3}}{3} - \frac{1}{3} i\right)^4$

 $\left(\frac{\sqrt{3}}{3} - \frac{1}{3} i\right)^4 =\answer{-\frac{8}{81}} + 
 \answer{-\frac{8\sqrt{3}}{81}} \, i$

\end{problem}

\begin{problem}
$\left(\frac{\sqrt{2}}{2} + \frac{\sqrt{2}}{2} i\right)^4$

\begin{solution}
 $\left(\frac{\sqrt{2}}{2} + \frac{\sqrt{2}}{2} i\right)^4=-1$
\end{solution}
\end{problem}

\begin{problem}
$(2+2i)^5$ 

$(2+2i)^5 = \answer{-128} + \answer{-128} i$

\end{problem}

\begin{problem}
$(\sqrt{3} - i)^{5}$ 

\begin{solution}
$(\sqrt{3} - i)^{5} =-16\sqrt{3} - 16i$
\end{solution}
\end{problem}

\begin{problem} \label{demoivrelast}
 $(1-i)^8$ 

 $(1-i)^8=\answer{16} + \answer{0} i$

\end{problem}

\end{question}

\begin{question}
In Exercises \ref{polarrootsfirst} - \ref{polarrootslast}, find the indicated complex roots.  Express your answers in polar form and then convert them into rectangular form.

\begin{problem} \label{polarrootsfirst}
the two square roots of $z = 4i$ 

\begin{solution}
Since $z=4i = 4\text{cis}\left(\frac{\pi}{2}\right)$ we have 

$w_{0} = 2\text{cis}\left(\frac{\pi}{4}\right) = \sqrt{2} +i\sqrt{2}$

$w_{1} = 2\text{cis}\left(\frac{5\pi}{4}\right) = -\sqrt{2} - i\sqrt{2}$


\end{solution}
\end{problem}

\begin{problem}
the two square roots of $z = -25i$

\begin{solution}
Since $z=-25i = 25\text{cis}\left(\frac{3\pi}{2}\right)$ we have 

$w_{0} = 5\text{cis}\left(\frac{3\pi}{4}\right) = -\frac{5\sqrt{2}}{2} +\frac{5\sqrt{2}}{2} i$

$w_{1} = 5\text{cis}\left(\frac{7\pi}{4}\right) = \frac{5\sqrt{2}}{2} - \frac{5\sqrt{2}}{2} i$

\end{solution}
\end{problem}

\begin{problem}
the two square roots of $z = 1 + i\sqrt{3}$

\begin{solution}

Since $z=1 + i\sqrt{3} = 2\text{cis}\left(\frac{\pi}{3}\right)$ we have 

$w_{0} = \sqrt{2}\text{cis}\left(\frac{\pi}{6}\right) = \frac{\sqrt{6}}{2} +\frac{\sqrt{2}}{2} i$

$w_{1} = \sqrt{2}\text{cis}\left(\frac{7\pi}{6}\right) = -\frac{\sqrt{6}}{2}-\frac{\sqrt{2}}{2} i$


\end{solution}
\end{problem}

\begin{problem}
the two square roots of $\frac{5}{2} - \frac{5\sqrt{3}}{2}i$

\begin{solution}

Since $z=\frac{5}{2} - \frac{5\sqrt{3}}{2}i = 5\text{cis}\left(\frac{5\pi}{3}\right)$ we have 

$w_{0} =\sqrt{5}\text{cis}\left(\frac{5\pi}{6}\right) = -\frac{\sqrt{15}}{2} + \frac{\sqrt{5}}{2}i$

$w_{1} = \sqrt{5}\text{cis}\left(\frac{11\pi}{6}\right) = \frac{\sqrt{15}}{2} - \frac{\sqrt{5}}{2}i$

\end{solution}

\end{problem}

\begin{problem}
the three cube roots of $z=64$

\begin{solution}

Since $z = 64 = 64\text{cis}\left(0\right)$ we have 

$w_{0} = 4\text{cis}\left(0\right) = 4$

$w_{1} =4\text{cis}\left(\frac{2\pi}{3}\right) = -2 + 2i\sqrt{3}$

$w_{2} = 4\text{cis}\left(\frac{4\pi}{3}\right) =  -2 - 2i\sqrt{3}$

\end{solution}
\end{problem}

\begin{problem}
the three cube roots of $z = -125$

\begin{solution}
Since $z = -125 = 125\text{cis}\left(\pi\right)$ we have 

$w_{0} = 5\text{cis}\left(\frac{\pi}{3}\right) = \frac{5}{2} + \frac{5\sqrt{3}}{2} i$

$w_{1} =5\text{cis}\left(\pi\right) = -5$

$w_{2} = 5\text{cis}\left(\frac{5\pi}{3}\right) = \frac{5}{2} - \frac{5\sqrt{3}}{2} i$
\end{solution}
\end{problem}

\begin{problem}
the three cube roots of $z = i$

\begin{solution}
Since $z = i = \text{cis}\left(\frac{\pi}{2}\right)$ we have 

$w_{0} = \text{cis}\left(\frac{\pi}{6}\right) = \frac{\sqrt{3}}{2} + \frac{1}{2}i$

$w_{1} = \text{cis}\left(\frac{5\pi}{6}\right) = -\frac{\sqrt{3}}{2} + \frac{1}{2}i$

$w_{2} = \text{cis}\left(\frac{3\pi}{2}\right) = -i$
\end{solution}
\end{problem}

\begin{problem}
the three cube roots of $z = -8i$

\begin{solution}
 Since $z = -8i = 8\text{cis}\left(\frac{3\pi}{2}\right)$ we have 

$w_{0} = 2\text{cis}\left(\frac{\pi}{2}\right) = 2i$

$w_{1} = 2\text{cis}\left(\frac{7\pi}{6}\right) = -\sqrt{3} -i$

$w_{2} = \text{cis}\left(\frac{11\pi}{6}\right) = \sqrt{3}-i$

\end{solution}
\end{problem}

\begin{problem}
the four fourth roots of $z=16$

\begin{solution}
Since $z=16 = 16\text{cis}\left(0 \right)$ we have 

$w_{0} =2\text{cis}\left(0\right) =2$

$w_{1} = 2\text{cis}\left(\frac{\pi}{2}\right) = 2i$

$w_{2} = 2\text{cis}\left(\pi\right) = -2$

$w_{3} = 2\text{cis}\left(\frac{3\pi}{2}\right) = -2i$

\end{solution}
\end{problem}

\begin{problem}
the four fourth roots of $z=-81$

\begin{solution}
Since $z=-81 = 81\text{cis}\left(\pi \right)$ we have 

$w_{0} =3\text{cis}\left(\frac{\pi}{4}\right) = \frac{3\sqrt{2}}{2} + \frac{3\sqrt{2}}{2}i$

$w_{1} = 3\text{cis}\left(\frac{3\pi}{4}\right) =-\frac{3\sqrt{2}}{2} + \frac{3\sqrt{2}}{2}i$

$w_{2} = 3\text{cis}\left(\frac{5\pi}{4}\right) =-\frac{3\sqrt{2}}{2} - \frac{3\sqrt{2}}{2}i$

$w_{3} = 3\text{cis}\left(\frac{7\pi}{4}\right) =\frac{3\sqrt{2}}{2} - \frac{3\sqrt{2}}{2}i$

\end{solution}
\end{problem}

\begin{problem}
the six sixth roots of $z = 64$

\begin{solution}
Since $z = 64 = 64\text{cis}(0)$ we have 

$w_{0} = 2\text{cis}(0) = 2$

$w_{1} = 2\text{cis}\left(\frac{\pi}{3}\right) = 1 + \sqrt{3}i$

$w_{2} = 2\text{cis}\left(\frac{2\pi}{3}\right) = -1 + \sqrt{3}i$

$w_{3} = 2\text{cis}\left(\pi\right) = -2$

$w_{4} = 2\text{cis}\left(-\frac{2\pi}{3}\right) = -1 - \sqrt{3}i$

$w_{5} = 2\text{cis}\left(-\frac{\pi}{3}\right) = 1 - \sqrt{3}i$

\end{solution}
\end{problem}

\begin{problem} \label{polarrootslast}
the six sixth roots of $z = -729$ 

\begin{solution}
Since $z = -729 = 729 \text{cis}(\pi)$ we have 

$w_{0} = 3\text{cis}\left(\frac{\pi}{6}\right) = \frac{3\sqrt{3}}{2} + \frac{3}{2}i$

$w_{1} = 3\text{cis}\left(\frac{\pi}{2}\right) = 3i$

$w_{2} = 3\text{cis}\left(\frac{5\pi}{6}\right) = -\frac{3\sqrt{3}}{2} + \frac{3}{2}i$

$w_{3} = 3\text{cis}\left(\frac{7\pi}{6}\right) =  -\frac{3\sqrt{3}}{2}-\frac{3}{2}i$

$w_{4} = 3\text{cis}\left(-\frac{3\pi}{2}\right) = -3i$

$w_{5} = 3\text{cis}\left(-\frac{11\pi}{6}\right) = \frac{3\sqrt{3}}{2} - \frac{3}{2}i$

\end{solution}

\end{problem}


\end{question}

\begin{problem}
Use the Sum and Difference Identities in Theorem \ref{circularsumdifference} or the Half Angle Identities in Theorem \ref{halfangle} to convert the three cube roots of $z=\sqrt{2} + i\sqrt{2}$ we found in  Example \ref{nthrootscomplexex}, number \ref{halfanglecuberoot} from polar form to rectangular form.

\begin{solution}

Note: In the answers for $w_{\text{\tiny$0$}}$ and $w_{\text{\tiny$2$}}$ the first rectangular form comes from applying the appropriate Sum or Difference Identity ($\frac{\pi}{12} = \frac{\pi}{3} - \frac{\pi}{4}$ and $\frac{17\pi}{12} = \frac{2\pi}{3} + \frac{3\pi}{4}$, respectively) and the second comes from using the Half-Angle Identities. 

$w_{0} = \sqrt[3]{2} \text{cis}\left(\frac{\pi}{12}\right) = \sqrt[3]{2}\left( \frac{\sqrt{6} + \sqrt{2}}{4} + i\left( \frac{\sqrt{6} - \sqrt{2}}{4} \right) \right) = \sqrt[3]{2}\left( \frac{\sqrt{2 + \sqrt{3}}}{2} + i\frac{\sqrt{2 - \sqrt{3}}}{2} \right)$ 
 
$w_{1} = \sqrt[3]{2} \text{cis}\left(\frac{3\pi}{4}\right) = \sqrt[3]{2} \left( -\frac{\sqrt{2}}{2} + \frac{\sqrt{2}}{2}i \right)$ 

$w_{2} = \sqrt[3]{2} \text{cis}\left(\frac{17\pi}{12}\right) = \sqrt[3]{2}\left( \frac{\sqrt{2} - \sqrt{6}}{4} + i\left( \frac{-\sqrt{2} - \sqrt{6}}{4} \right) \right) = \sqrt[3]{2}\left( \frac{\sqrt{2 - \sqrt{3}}}{2} + i\frac{\sqrt{2 + \sqrt{3}}}{2} \right)$

\end{solution}
\end{problem}

\begin{problem}
Use a calculator to approximate the rectangular form of the five fifth roots of $1$ we found in Example \ref{nthrootscomplexex}, number \ref{calculatorfifthroot}.

\begin{solution}

 $w_{0} = \text{cis}(0) = 1$

$w_{1} = \text{cis}\left(\frac{2\pi}{5}\right) \approx 0.309 + 0.951i$

$w_{2} = \text{cis}\left(\frac{4\pi}{5}\right) \approx -0.809 + 0.588i$

$w_{3} = \text{cis}\left(\frac{6\pi}{5}\right) \approx -0.809 - 0.588i$

$w_{4} = \text{cis}\left(\frac{8\pi}{5}\right) \approx 0.309 - 0.951i$

\end{solution}
\end{problem}

\begin{problem}
According to Theorem \ref{realfactorization} in Section \ref{ComplexZeros}, the polynomial $p(x) = x^{4} + 4$ can be factored into the product linear and irreducible quadratic factors.  In Exercise \ref{factorpolywithnonlinear} in Section \ref{NonLinearEquations}, we showed you how to factor this polynomial into the product of two irreducible quadratic factors using a system of non-linear equations.  Now that we can compute the complex fourth roots of $-4$ directly using  Theorem \ref{nthrootscomplexthm}, we can  apply the Complex Factorization Theorem, Theorem \ref{complexfactorization}, to obtain the linear factorization $p(x) = (x - (1 + i))(x - (1 - i))(x - (-1 + i))(x - (-1 - i))$.  By multiplying the first two factors together and then the second two factors together, thus pairing up the complex conjugate pairs of zeros Theorem \ref{conjugatepairsthm} told us we'd get, we have that $p(x) = (x^{2} - 2x + 2)(x^{2} + 2x + 2)$.  Use the 12 complex $12^{\text{th}}$ roots of 4096 to factor $p(x) = x^{12} - 4096$ into a product of linear and irreducible quadratic factors.

\begin{solution}
$p(x) = x^{12} - 4096 = (x - 2)(x + 2)(x^{2} + 4)(x^{2} - 2x + 4)(x^{2} + 2x + 4)(x^{2} - 2\sqrt{3}x + 4)(x^{2} + 2\sqrt{3} + 4)$
\end{solution}
\end{problem}


\begin{problem}
Use Exercise \ref{triangleineqforvectorsexercise} from Section \ref{TheDotProduct} to show the the Triangle Inequality $|z + w| \leq |z| + |w|$ holds for all complex numbers $z$ and $w$ as well.  Identify the complex number $z = a + bi$ with the vector $u = \langle a, b \rangle$ and identify the complex number $w = c + di$ with the vector $v = \langle c, d \rangle$ and just follow your nose!
\end{problem}


\begin{problem}
Complete the proof of Theorem \ref{modprops} by showing that if $w \neq 0$ than  $\left| \frac{1}{w}\right| = \frac{1}{|w|}$.
\end{problem}

\begin{problem}
Recall from Section \ref{ComplexZeros} that given a complex number $z = a+bi$ its complex conjugate, denoted $\overline{z}$, is given by $\overline{z} = a - bi$.

\begin{enumerate}

\item Prove that $\left| \overline{z} \right| = |z|$.

\item Prove that $|z| = \sqrt{z \overline{z}}$

\item Show that $\text{Re}(z) = \frac{z + \overline{z}}{2}$ and $\text{Im}(z) = \frac{z - \overline{z}}{2i}$

\item Show that if $\theta \in \text{arg}(z)$ then $-\theta \in \text{arg}\left(\overline{z}\right)$. Interpret this result geometrically.

 \item Is it always true that $\text{Arg}\left(\overline{z}\right) = -\text{Arg}(z)$?
\end{enumerate}

\end{problem}

\begin{problem}
Given a natural number $n \geq 2$, the $n$ complex $n^{\text{th}}$ roots of  $z = 1$ are called the \textbf{\boldmath $n^{\mbox{\textbf{\scriptsize th}}}$ Roots of Unity}. \index{$n^{\textrm{th}}$ Roots of Unity} \index{complex number ! $n^{\textrm{th}}$ Roots of Unity} \index{Roots of Unity} In the following exercises, assume that $n$ is a fixed, but arbitrary, natural number such that $n \geq 2$.

\begin{enumerate}

\item Show that $w = 1$ is an $n^{\text{th}}$ root of unity.

\item Show that if both $w_{\text{\tiny$j$}}$ and $w_{\text{\tiny$k$}}$ are $n^{\text{th}}$ roots of unity then so is their product $w_{\text{\tiny$j$}}w_{\text{\tiny$k$}}$.

\item Show that if $w_{\text{\tiny$j$}}$ is an $n^{\text{th}}$ root of unity then there is an $n^{\text{th}}$ root of unity $w_{\text{\tiny$j'$}}$  so that  $w_{\text{\tiny$j$}}w_{\text{\tiny$j'$}} = 1$. 



 HINT: If $w_{\text{\tiny$j$}} = \text{cis}(\theta)$ let $w_{\text{\tiny$j'$}} = \text{cis}(2\pi - \theta)$. Show $w_{\text{\tiny$j'$}} = \text{cis}(2\pi - \theta)$ is indeed an $n^{\text{th}}$ root of unity.

\end{enumerate}

\end{problem}

\begin{problem} \label{eulerformulaexercise} 
Another way to express the polar form of a complex number is to use the exponential function.  For real numbers $t$, \href{http://en.wikipedia.org/wiki/Leonhard_Euler}{\underline{Euler}}'s Formula defines $e^{it} = \cos(t) + i \sin(t)$.  

\begin{enumerate}

\item Use Theorem \ref{prodquotpolarcomplex} to show that:

\begin{enumerate}

\item  $e^{ix} e^{iy} = e^{i(x+y)}$ for all real numbers $x$ and $y$.

\item  $\left(e^{ix}\right)^{n} = e^{i(nx)}$ for any real number $x$ and any natural number $n$.

\item $\frac{e^{ix}}{e^{iy}} = e^{i(x-y)}$ for all real numbers $x$ and $y$.

\end{enumerate}



\item If $z = r\text{cis}(\theta)$ is the polar form of $z$, show that $z = re^{it}$ where $\theta = t$ radians.

\item Show that $e^{i\pi} + 1 = 0$.  (This famous equation relates the five most important constants in all of Mathematics with the three most fundamental operations in Mathematics.)

\item  \label{expformcosandsin} Show that $\cos(t) = \frac{e^{it} + e^{-it}}{2}$ and that $\sin(t) = \frac{e^{it} - e^{-it}}{2i}$ for all real numbers $t$. 

\end{enumerate}

\end{problem}

\end{document}

